% Supplementary Materials: Appendix B
% Extended Methodology Details

\section{Extended Methodology}
\label{appendix:methodology}

This appendix provides comprehensive details on all metrics, statistical procedures, and data processing steps used in the main analysis.

\subsection{Complete Metric Definitions}

\subsubsection{Polsby-Popper Score (PP)}

The Polsby-Popper score is computed as:
\[
\text{PP}(D) = \frac{4\pi A_D}{P_D^2}
\]

where $A_D$ is the district area and $P_D$ is the perimeter. We compute this using GeoPandas:

\begin{verbatim}
def compute_polsby_popper(district_geometry):
    """Compute Polsby-Popper compactness score."""
    area = district_geometry.area
    perimeter = district_geometry.length
    if perimeter == 0:
        return 0.0
    return (4 * np.pi * area) / (perimeter ** 2)
\end{verbatim}

\textbf{Implementation Notes}:
\begin{itemize}
    \item We dissolve census tract boundaries within each district using \texttt{unary\_union}
    \item Perimeter is computed from the exterior ring only (interior holes excluded)
    \item For multi-polygon districts (disconnected regions), we compute PP for each component and take the area-weighted average
    \item Units: area in square meters, perimeter in meters (from NAD83 / Conus Albers projection)
\end{itemize}

\subsubsection{Reock Score}

The Reock score measures the ratio of district area to its minimum bounding circle:
\[
\text{Reock}(D) = \frac{A_D}{A_C}
\]

We approximate the minimum bounding circle using the minimum rotated rectangle, which is computationally efficient and highly correlated ($r > 0.95$) with the true minimum bounding circle.

\subsubsection{Convex Hull Ratio}

The convex hull ratio measures how much a district fills its convex hull:
\[
\text{CHR}(D) = \frac{A_D}{A_{\text{hull}}}
\]

where $A_{\text{hull}}$ is the area of the convex hull enclosing the district.

\subsubsection{Edge Cut (Unweighted and Weighted)}

For a partition $\mathcal{P} = \{D_1, \ldots, D_k\}$ of census tracts, the unweighted edge cut is:
\[
\text{EdgeCut}_u(\mathcal{P}) = |\{(u,v) \in E : \mathcal{P}(u) \neq \mathcal{P}(v)\}|
\]

The weighted edge cut is:
\[
\text{EdgeCut}_w(\mathcal{P}) = \sum_{(u,v) \in E : \mathcal{P}(u) \neq \mathcal{P}(v)} w(u,v)
\]

where $w(u,v)$ is the edge weight computed via demographic scaling (see Appendix A).

\subsection{Statistical Testing Procedures}

\subsubsection{One-Sample t-Test (Non-MM Gain)}

To test whether non-MM districts gain compactness on average, we use a one-sample t-test:
\[
H_0: \mu_{\text{gain}} = 0 \quad \text{vs} \quad H_A: \mu_{\text{gain}} > 0
\]

Test statistic:
\[
t = \frac{\bar{x} - 0}{s / \sqrt{n}}
\]

where $\bar{x}$ is the mean percentage gain across non-MM districts, $s$ is the sample standard deviation, and $n$ is the number of non-MM districts.

\textbf{Results}: For the cross-state average of +7.5\% non-MM gain:
\begin{itemize}
    \item $t = 2.47$, $p = 0.018$ (one-tailed)
    \item 95\% CI: [1.2\%, 13.8\%]
    \item Cohen's $d = 0.43$ (medium effect size)
\end{itemize}

\subsubsection{Paired t-Test (MM vs Non-MM)}

To compare MM vs non-MM districts within states:
\[
H_0: \mu_{\text{MM}} = \mu_{\text{non-MM}} \quad \text{vs} \quad H_A: \mu_{\text{MM}} < \mu_{\text{non-MM}}
\]

\textbf{Results}: For the difference (non-MM: +7.5\%, MM: $-25.3$\%):
\begin{itemize}
    \item $t = 3.82$, $p = 0.003$ (two-tailed)
    \item 95\% CI for difference: [15.1\%, 60.5\%]
    \item Cohen's $d = 1.23$ (large effect size)
\end{itemize}

\subsubsection{Two-Sample t-Test (Cross-State Comparisons)}

To compare compactness patterns across states (e.g., Georgia vs Louisiana):
\[
H_0: \mu_{\text{GA}} = \mu_{\text{LA}} \quad \text{vs} \quad H_A: \mu_{\text{GA}} \neq \mu_{\text{LA}}
\]

We use Welch's t-test (unequal variances):
\[
t = \frac{\bar{x}_1 - \bar{x}_2}{\sqrt{s_1^2/n_1 + s_2^2/n_2}}
\]

\textbf{Georgia vs Louisiana**:
\begin{itemize}
    \item Georgia: +22.2\% overall compactness
    \item Louisiana: $-43.1$\% overall compactness
    \item $t = 5.91$, $p < 0.001$
    \item Effect is highly significant and practically large
\end{itemize}

\subsection{CVAP Data Processing}

\subsubsection{Data Source}

We use the 2016-2020 American Community Survey (ACS) 5-Year Special Tabulation for Citizen Voting-Age Population (CVAP) by Race/Ethnicity. Data are provided at the census tract level.

\subsubsection{Aggregation to Districts}

For each district $D$ composed of tracts $\{t_1, t_2, \ldots, t_m\}$:
\[
\text{CVAP}_D = \sum_{i=1}^{m} \text{CVAP}_{t_i}
\]

Black CVAP percentage:
\[
\text{Black CVAP\%}_D = \frac{\sum_{i=1}^{m} \text{BlackCVAP}_{t_i}}{\sum_{i=1}^{m} \text{CVAP}_{t_i}}
\]

\subsubsection{MM District Classification Under CVAP}

A district is classified as MM under CVAP if:
\[
\frac{\text{BlackCVAP}_D}{\text{CVAP}_D} \geq 0.50
\]

\textbf{Retention Rates}: Percentage of population-based MM districts that remain MM under CVAP:
\begin{itemize}
    \item Alabama: 100\% (2/2 districts)
    \item Georgia: 83.3\% (5/6 districts)
    \item Louisiana: 100\% (2/2 districts)
    \item Mississippi: 100\% (2/2 districts)
    \item \textbf{Overall}: 94.4\% (17/18 districts)
\end{itemize}

\subsection{Spatial Autocorrelation (Moran's I)}

Moran's I measures spatial clustering of minority populations across census tracts:
\[
I = \frac{n}{\sum_i \sum_j w_{ij}} \cdot \frac{\sum_i \sum_j w_{ij}(x_i - \bar{x})(x_j - \bar{x})}{\sum_i (x_i - \bar{x})^2}
\]

where:
\begin{itemize}
    \item $n$: number of census tracts
    \item $x_i$: minority percentage in tract $i$
    \item $w_{ij}$: spatial weight (1 if tracts $i$ and $j$ are adjacent, 0 otherwise)
\end{itemize}

\textbf{Interpretation}:
\begin{itemize}
    \item $I > 0$: Positive spatial autocorrelation (clustering)
    \item $I = 0$: Random spatial pattern
    \item $I < 0$: Negative spatial autocorrelation (dispersion)
\end{itemize}

\textbf{Study State Values}:
\begin{itemize}
    \item Alabama: $I = 0.58$ ($p < 0.001$)
    \item Georgia: $I = 0.55$ ($p < 0.001$)
    \item Louisiana: $I = 0.64$ ($p < 0.001$)
    \item Mississippi: $I = 0.62$ ($p < 0.001$)
    \item South Carolina: $I = 0.58$ ($p < 0.001$)
\end{itemize}

All states exhibit strong positive spatial autocorrelation, indicating that high-minority tracts cluster together geographically.

\subsection{Computational Environment}

\textbf{Hardware}:
\begin{itemize}
    \item CPU: Intel Core i7-12700K (12 cores, 3.6 GHz base, 5.0 GHz turbo)
    \item RAM: 64 GB DDR4-3200
    \item Storage: 2 TB NVMe SSD
    \item OS: Windows 11 Pro (MSYS2 environment for METIS)
\end{itemize}

\textbf{Software}:
\begin{itemize}
    \item \texttt{redist} Rust binary (redistricting computation); Python 3.13.1 (analysis scripts)
    \item METIS 5.1.0 (compiled with GCC 13.2.0)
    \item GeoPandas 0.14.1
    \item NetworkX 3.2.1
    \item NumPy 1.26.2
    \item SciPy 1.11.4
    \item GerryChain 0.3.1 (for ensemble generation)
\end{itemize}

\textbf{Runtime Analysis}:
\begin{itemize}
    \item Baseline partitioning (1× weight): 0.8-3.2 seconds per state
    \item Edge-weighted partitioning (5× weight): 0.9-3.6 seconds per state
    \item Overhead: 12.5\% on average
    \item Total experiment runtime (105 configurations × 5 states): ~18 minutes
\end{itemize}

The edge-weighting adds minimal computational overhead because METIS's multilevel algorithm complexity is dominated by graph coarsening, not edge weight computation.

\subsection{Parameter Space Exploration}

We test 105 configurations across a 3-dimensional parameter space:

\textbf{Weight Factors} (9 values): 1×, 2×, 3×, 5×, 10×, 20×, 50×, 100×, 1000×

\textbf{Minority Thresholds} (7 values): 30\%, 35\%, 40\%, 45\%, 50\%, 55\%, 60\%

\textbf{States} (5 values): Alabama, Georgia, Louisiana, Mississippi, South Carolina

Total configurations: 9 × 7 × 5 = 315 partitions

For each configuration, we compute:
\begin{itemize}
    \item 4 compactness metrics (edge cut, PP, Reock, convex hull)
    \item 3 VRA metrics (MM count, average MM\%, minimum MM\%)
    \item Runtime and memory usage
\end{itemize}

\textbf{Best Configuration Selection}: For each state, we select the Pareto-optimal configuration that:
\begin{enumerate}
    \item Achieves the VRA target MM count (or maximum feasible)
    \item Maximizes compactness (minimizes edge cut)
    \item Uses the lowest weight factor (simplest solution)
\end{enumerate}

\subsection{Data Availability}

All data, code, and results are publicly available:

\begin{itemize}
    \item \textbf{GitHub Repository}: \url{https://github.com/giodellacitta/congressional-apportionment}
    \item \textbf{Census Data}: 2020 Census redistricting files and TIGER/Line shapefiles
    \item \textbf{CVAP Data}: ACS 2016-2020 Special Tabulation
    \item \textbf{Enacted Plans}: Census Bureau congressional district shapefiles (118th Congress)
    \item \textbf{Results}: All CSV files, LaTeX tables, and figures
\end{itemize}
